% =============================================================
% §3 Problem Formulation
% =============================================================
\section{Problem Formulation}
\label{sec:formulation}

We now formalize the cloud-edge LLM inference scheduling problem.
\S\ref{sec:sys-model} introduces notation for the deployment
topology and workload. \S\ref{sec:physics-constraints} expresses
the five AIGC physics as constraints and cost functions on
scheduling decisions. \S\ref{sec:sched-problem} states our
scheduling problem as a multi-objective optimization.
\S\ref{sec:mdp} casts it as a Markov decision process (MDP) for
reinforcement-learning treatment. Table~\ref{tab:notation}
summarizes the notation used throughout this section.

\begin{table}[t]
  \centering
  \caption{Notation used in \S\ref{sec:formulation}.}
  \label{tab:notation}
  \small
  \setlength{\tabcolsep}{4pt}
  \begin{tabular}{@{}l@{\hspace{8pt}}p{0.60\columnwidth}@{}}
    \toprule
    Symbol & Meaning \\
    \midrule
    \multicolumn{2}{@{}l}{\emph{Topology and servers}} \\
    $\mathcal{S} = \{s_0, \ldots, s_N\}$
      & server pool: cloud $s_0$, edge $s_1, \ldots, s_N$ \\
    $C_s$, $V_s$, $B_s$
      & compute (TFLOPS), VRAM (GB), bandwidth (Mbps) of $s$ \\
    $P^{\text{idle}}_s$, $P^{\max}_s$
      & idle and peak power of $s$ (W) \\
    $u_s(t)$, $P_s(t)$
      & utilization and instantaneous power of $s$ \\
    $\tau(s, s', x)$
      & transfer time of $x$~GB from $s$ to $s'$ \\
    \addlinespace
    \multicolumn{2}{@{}l}{\emph{Model catalog}} \\
    $\mathcal{M} = \{m_1, \ldots, m_K\}$
      & catalog of LLM variants \\
    $W_m$, $C^{\text{cold}}_m$
      & weight footprint (GB) and cold-load time (s) of $m$ \\
    $\kappa_m$, $\phi_m$, $B^{\max}_m$
      & per-token KV size, decode floor, max batch of $m$ \\
    $L_s(t)$
      & set of models loaded on $s$ at time $t$ \\
    \addlinespace
    \multicolumn{2}{@{}l}{\emph{Workload}} \\
    $\mathcal{R}$, $\lambda$
      & request stream and Poisson arrival rate (req/s) \\
    $r_i = (m_i, L^p_i, L^o_i, a_i)$
      & target model, prompt/output length, arrival time \\
    $t^{\text{pre}}_i$, $t^{\text{dec}}_i$
      & prefill and decode task of request $r_i$ \\
    $w_t$, $n_t$
      & workload (TFLOPS) and token count of task $t$ \\
    \addlinespace
    \multicolumn{2}{@{}l}{\emph{Execution and objectives}} \\
    $\delta^{\text{cold}}$, $\delta^{\text{kv}}$
      & cold-load and KV-migration delay \\
    $b_{s,m,k}(t)$, $\beta_k$
      & batch occupancy; phase-$k$ overhead \\
    $T^{\text{solo}}$, $T^{\text{exec}}_{\text{eff}}$
      & solo / batched execution time \\
    $e^{\text{pre}}_i$; $s^{\text{dec}}_i$, $e^{\text{dec}}_i$
      & prefill completion; decode start and completion \\
    $\text{TTFT}_i$, $\text{TPOT}_i$
      & time-to-first-token, time-per-output-token \\
    $T^*_{\text{TTFT}}$, $T^*_{\text{TPOT}}$
      & \slo{} thresholds on TTFT and TPOT \\
    $\text{SLO}(\pi)$, $E_{\text{tok}}(\pi)$
      & \slo{} attainment and energy per token under $\pi$ \\
    $\pi$; $\mathbf{s}_t$, $a_t$, $r_t$
      & scheduling policy; MDP state, action, reward \\
    \bottomrule
  \end{tabular}
\end{table}

\subsection{System Model}
\label{sec:sys-model}

\textbf{Deployment topology.} A heterogeneous compute pool
$\mathcal{S} = \{s_0, s_1, \ldots, s_N\}$ where $s_0$ is a cloud
server and $s_1, \ldots, s_N$ are edge servers. Each server $s$
has fixed attributes: compute capacity $C_s$ (TFLOPS), VRAM
capacity $V_s$ (GB), uplink bandwidth $B_s$ (Mbps), idle and
peak power $P^{\text{idle}}_s$ and $P^{\max}_s$ (W). A network
function $\tau(s, s', x)$ returns the transfer time of $x$~GB
of data between servers $s$ and $s'$.

\textbf{Model catalog.} A set
$\mathcal{M} = \{m_1, \ldots, m_K\}$ of LLM variants. Each model
$m$ has weight footprint $W_m$ (GB), cold-load time
$C^{\text{cold}}_m$ (s), per-token KV-cache size $\kappa_m$
(MB/token), per-token decode floor latency $\phi_m$ (s/token),
and maximum batch size $B^{\max}_m$.

\textbf{Request workload.} A stream of inference requests
$\mathcal{R} = \{r_1, r_2, \ldots\}$ arrives over time according
to a Poisson process at rate $\lambda$. Each request
$r_i = (m_i, L^p_i, L^o_i, a_i)$ has a target model
$m_i \in \mathcal{M}$, prompt length $L^p_i$, expected output
length $L^o_i$, and arrival time $a_i$. The target model is
\emph{part of the request}, mirroring production APIs where
each call names its model endpoint; model selection happens
upstream---the scheduler decides \emph{where} each request
runs, never \emph{which} model serves it.

\textbf{Task decomposition.} Each request $r_i$ decomposes into
two tasks $(t^{\text{pre}}_i, t^{\text{dec}}_i)$ with a strict
dependency edge $t^{\text{pre}}_i \to t^{\text{dec}}_i$. The
prefill task carries workload
$w^{\text{pre}}_i \propto L^p_i$, the decode task carries
$w^{\text{dec}}_i \propto L^o_i$, and both reference the same
model $m_i$ and request id $i$. The inference workload is
therefore a forest of pairwise-independent two-node DAGs.

\subsection{AIGC Physics as Constraints}
\label{sec:physics-constraints}

The five AIGC physics from \S\ref{sec:intro} manifest as the
following constraints and cost terms governing how tasks
execute on servers.

\textbf{(P1) Model weight residency.} Server $s$ holds a state
$L_s(t) \subseteq \mathcal{M}$ representing the set of models
currently loaded in its VRAM at time $t$. A task targeting
model $m$ admitted to $s$ pays a cold-load penalty
$C^{\text{cold}}_m$ if $m \notin L_s(t)$. If loading $m$ would
exceed $V_s$, the simulator evicts unpinned models in LRU
order. Formally, the admission delay is:
\begin{equation}
  \delta^{\text{cold}}_{s,m,t} = C^{\text{cold}}_m \cdot
    \mathbb{1}[m \notin L_s(t)].
\end{equation}

\textbf{(P2) KV-cache locality.} For each decode task
$t^{\text{dec}}_i$ scheduled to server $s$ when its prefill
$t^{\text{pre}}_i$ ran on $s'$, an additional KV-migration
cost applies:
\begin{equation}
  \delta^{\text{kv}}_i = \tau\!\left(s', s,
    L^p_i \cdot \kappa_{m_i} / 1024\right).
\end{equation}
When $s = s'$, this cost is zero (KV cache is local).

\textbf{(P3) Continuous batching.} Let $b_{s,m,k}(t)$ denote
the number of running same-model same-kind tasks on server $s$
for model $m$, phase $k \in \{\text{pre}, \text{dec}\}$. A new
task admitted at time $t$ to $(s, m, k)$ experiences effective
execution time
\begin{equation}
  T^{\text{exec}}_{\text{eff}}
    = T^{\text{solo}} \cdot
      \bigl(1 + b_{s,m,k}(t) \cdot \beta_k\bigr),
\end{equation}
where $\beta_{\text{pre}} = 0.30$ and
$\beta_{\text{dec}} = 0.05$ are phase-specific overhead
factors, calibrated to vLLM measurements
(\S\ref{sec:simulator}).
Admission requires $b_{s,m,k}(t) < B^{\max}_m$ (batch slot not
full).

\textbf{(P4) Memory-bandwidth floor.} The solo execution time
of a phase task is the maximum of its compute-bound time and
its memory-bound floor:
\begin{equation}
  T^{\text{solo}}
    = \max\!\left(\frac{w_t}{C_s},\;
                  \phi_{m_t} \cdot n_t\right),
\end{equation}
where $w_t$ is the task's workload (TFLOPS) and $n_t$ is its
token count (input for prefill, output for decode).

\textbf{(P5) Two-phase dependency.} The decode task
$t^{\text{dec}}_i$ cannot become ready until
$t^{\text{pre}}_i$ has fully completed. Combined with (P1) and
(P2), this means a decode dispatched to $s \neq s'$ sees no
benefit from prefill's KV cache---either it pays migration
cost or re-runs prefill.

\subsection{Scheduling Problem}
\label{sec:sched-problem}

A scheduling policy $\pi$ produces, for each ready task at each
time step, an assignment $\pi: t \mapsto s \in \mathcal{S}$.
Given a request stream $\mathcal{R}$, $\pi$ induces:

\begin{itemize}
\item Per-request \textbf{TTFT} (time-to-first-token):
  $\text{TTFT}_i = e^{\text{pre}}_i - a_i$, where
  $e^{\text{pre}}_i$ is the wall-clock completion time of the
  prefill task.

\item Per-request \textbf{TPOT} (time-per-output-token):
  $\text{TPOT}_i = (e^{\text{dec}}_i - s^{\text{dec}}_i) / L^o_i$.

\item \textbf{\slo{} attainment}:
  \begin{equation}
    \text{SLO}(\pi) = \frac{1}{|\mathcal{R}|} \sum_i
      \mathbb{1}\!\bigl[\text{TTFT}_i \le T^*_{\text{TTFT}}
        \land \text{TPOT}_i \le T^*_{\text{TPOT}}\bigr].
  \end{equation}

\item \textbf{Energy per token}:
  \begin{equation}
    E_{\text{tok}}(\pi) =
      \frac{\sum_s \int_0^T P_s(t)\,dt}{\sum_i L^o_i},
  \end{equation}
  where instantaneous power
  $P_s(t) = P^{\text{idle}}_s + (P^{\max}_s - P^{\text{idle}}_s)
  \cdot u_s(t)$ depends on per-server compute utilization
  $u_s(t)$.
\end{itemize}

Our scheduling problem is the \textbf{bi-objective program}:
\begin{equation}
  \max_\pi\; \text{SLO}(\pi)
  \qquad
  \min_\pi\; E_{\text{tok}}(\pi),
\end{equation}
subject to admission feasibility (compute, VRAM, batch-slot)
at every step. Because the two objectives can be in tension
(e.g., always-cloud maximizes per-request compute but
overshoots energy budget), we evaluate policies on the
$(\text{SLO}, E_{\text{tok}})$ plane and seek \textbf{Pareto
dominance}: $\pi$ dominates $\pi'$ if $\pi$ is no worse on both
axes and strictly better on at least one. These two objectives
operationalize the latency and cost motivations from
\S\ref{sec:intro}; privacy, when required, enters as a
feasibility constraint on the admissible server set rather
than as an objective.

\subsection{MDP Formulation for RL}
\label{sec:mdp}

We treat per-task dispatch as a sequential decision process.
At each ready-task arrival, the agent observes a state
$\mathbf{s}_t$ summarizing
the task's properties and all servers' utilization,
loaded-model sets, and batch occupancies (\S\ref{sec:rl}),
emits an action
$a_t \in \mathcal{S}$ choosing a server, and receives a scalar
reward $r_t$: a weighted sum of seven physical and AIGC-aware
terms (\S\ref{sec:rl}) providing
\textbf{dense per-decision} feedback rather than
episode-end \slo{}/energy outcomes---essential
given the long episode horizons ($\sim$100 dispatches per
run).

The MDP is \textbf{partially observable} in principle, but
the observable features capture all decision-relevant
information under the simulator's dynamics; it is
non-stationary within an episode (utilization
evolves) and stationary across
episodes (fixed workload distribution). We solve it with PPO
with action masking, pretraining, and generalized advantage
estimation (GAE), as detailed in
\S\ref{sec:rl}.

\textbf{Why RL.} A natural alternative to RL is solving the
integer-linear program (ILP) implied by \S\ref{sec:sched-problem}
directly. We rejected this because (i) AIGC physics (P1)(P3)
introduce \emph{state-dependent} costs (cold-load and batch
overhead depend on server's running set, which changes within
an episode), breaking the static-cost ILP assumption; and
(ii) the combinatorial action space at each step
($|\mathcal{S}|^{|\mathcal{R}|}$) is too large for
branch-and-bound at realistic workload sizes.
